\textbf{Phần 1. Trắc nghiệm lựa chọn}
\setcounterpageref{ex}{0}

\begin{ex} %[5D7V43]
Đi siêu thị, mẹ mua $2{,}5\text{ kg}$ thịt với giá $130\,000$ đồng/kg, $3\text{ kg}$ cá với giá $65\,000$ đồng/kg và rau củ quả hết $45\,000$ đồng. Sử dụng máy tính cầm tay tính tổng số tiền mẹ phải thanh toán là
\choice
{$520\,000\text{ đồng}$}
{\True $565\,000\text{ đồng}$}
{$585\,000\text{ đồng}$}
{$615\,000\text{ đồng}$}
\loigiai{
Số tiền mua thịt: $2{,}5 \times 130\,000 = 325\,000\text{ (đồng)}$.\\
Số tiền mua cá: $3 \times 65\,000 = 195\,000\text{ (đồng)}$.\\
Tổng số tiền cần trả là:
$325\,000 + 195\,000 + 45\,000 = 565\,000\text{ (đồng)}$.\\
\textbf{Đáp án của câu hỏi là B}
}
\end{ex}

\begin{ex} %[5D7V43]
Hóa đơn tiền điện tháng này của một hộ gia đình là $750\,000$ đồng. Thuế giá trị gia tăng (VAT) được tính thêm là $8\%$. Tổng số tiền gia đình đó phải thanh toán là
\choice
{$790\,000\text{ đồng}$}
{$800\,000\text{ đồng}$}
{\True $810\,000\text{ đồng}$}
{$820\,000\text{ đồng}$}
\loigiai{
Tiền thuế VAT là:
$750\,000 \times 8\% = 60\,000\text{ (đồng)}$.\\
Tổng số tiền phải thanh toán là:
$750\,000 + 60\,000 = 810\,000\text{ (đồng)}$.\\
\textbf{Đáp án của câu hỏi là C}
}
\end{ex}

\begin{ex} %[5D7V43]
Một mảnh đất hình thang có đáy lớn $42{,}5\text{ m}$, đáy bé $27{,}5\text{ m}$ và chiều cao $18\text{ m}$. Sử dụng máy tính cầm tay tính diện tích mảnh đất đó là
\choice
{$580\text{ m}^2$}
{\True $630\text{ m}^2$}
{$1\,260\text{ m}^2$}
{$720\text{ m}^2$}
\loigiai{
Diện tích mảnh đất là:
$(42{,}5 + 27{,}5) \times 18 : 2 = 70 \times 9 = 630\text{ (m}^2\text{)}$.\\
\textbf{Đáp án của câu hỏi là B}
}
\end{ex}

\begin{ex} %[5D7V43]
Bác Ba mua $15$ bao xi măng với giá $85\,000$ đồng/bao và $2$ chuyến xe cát với giá $450\,000$ đồng/chuyến. Bác Ba đưa cho người bán hàng $2\,500\,000$ đồng. Số tiền thừa bác Ba nhận lại là
\choice
{$300\,000\text{ đồng}$}
{\True $325\,000\text{ đồng}$}
{$350\,000\text{ đồng}$}
{$375\,000\text{ đồng}$}
\loigiai{
Tổng số tiền mua hàng là:
$15 \times 85\,000 + 2 \times 450\,000 = 1\,275\,000 + 900\,000 = 2\,175\,000\text{ (đồng)}$.\\
Số tiền thừa nhận lại là:
$2\,500\,000 - 2\,175\,000 = 325\,000\text{ (đồng)}$.\\
\textbf{Đáp án của câu hỏi là B}
}
\end{ex}

\begin{ex} %[5D7V43]
Bác An gửi tiết kiệm $80\,000\,000$ đồng với lãi suất $0{,}45\%$ mỗi tháng. Sau $3$ tháng, số tiền lãi bác An nhận được là
\choice
{$360\,000\text{ đồng}$}
{$720\,000\text{ đồng}$}
{\True $1\,080\,000\text{ đồng}$}
{$1\,200\,000\text{ đồng}$}
\loigiai{
Tiền lãi một tháng là:
$80\,000\,000 \times 0{,}45\% = 360\,000\text{ (đồng)}$.\\
Sau 3 tháng số tiền lãi nhận được là:
$360\,000 \times 3 = 1\,080\,000\text{ (đồng)}$.\\
\textbf{Đáp án của câu hỏi là C}
}
\end{ex}

\begin{ex} %[5D7V43]
Một chiếc xe máy tiêu thụ trung bình $2{,}4\text{ l}$ xăng khi đi quãng đường $100\text{ km}$. Giá tiền $1\text{ l}$ xăng là $23\,500$ đồng. Chi phí tiền xăng để xe máy đi hết quãng đường $250\text{ km}$ là
\choice
{$135\,000\text{ đồng}$}
{\True $141\,000\text{ đồng}$}
{$145\,000\text{ đồng}$}
{$152\,000\text{ đồng}$}
\loigiai{
Số lít xăng tiêu thụ cho $250\text{ km}$ là:
$2{,}4 \times (250 : 100) = 6\text{ (l)}$.\\
Chi phí tiền xăng là:
$6 \times 23\,500 = 141\,000\text{ (đồng)}$.\\
\textbf{Đáp án của câu hỏi là B}
}
\end{ex}

\begin{ex} %[5D7V43]
Một siêu thị điện máy niêm yết giá bán một chiếc máy giặt là $7\,500\,000$ đồng và có chương trình giảm giá $12\%$. Nếu khách thanh toán bằng ví điện tử thì được giảm thêm $100\,000$ đồng. Số tiền thực tế khách cần thanh toán là
\choice
{$6\,400\,000\text{ đồng}$}
{\True $6\,500\,000\text{ đồng}$}
{$6\,600\,000\text{ đồng}$}
{$6\,700\,000\text{ đồng}$}
\loigiai{
Giá sau khi giảm $12\%$ là:
$7\,500\,000 \times (100\% - 12\%) = 6\,600\,000\text{ (đồng)}$.\\
Số tiền thanh toán cuối cùng là:
$6\,600\,000 - 100\,000 = 6\,500\,000\text{ (đồng)}$.\\
\textbf{Đáp án của câu hỏi là B}
}
\end{ex}

\begin{ex} %[5D7V43]
Một thửa ruộng hình chữ nhật có chiều dài $64\text{ m}$ và chiều rộng $25\text{ m}$. Cứ mỗi $100\text{ m}^2$ người ta thu hoạch được $65\text{ kg}$ lúa. Khối lượng lúa thu hoạch được trên cả thửa ruộng là
\choice
{$980\text{ kg}$}
{$1\,000\text{ kg}$}
{\True $1\,040\text{ kg}$}
{$1\,080\text{ kg}$}
\loigiai{
Diện tích thửa ruộng là:
$64 \times 25 = 1\,600\text{ (m}^2\text{)}$.\\
Khối lượng lúa thu hoạch được là:
$1\,600 : 100 \times 65 = 1\,040\text{ (kg)}$.\\
\textbf{Đáp án của câu hỏi là C}
}
\end{ex}

\begin{ex} %[5D7V43]
Một bể bơi hình hộp chữ nhật có lòng trong dài $25\text{ m}$, rộng $12\text{ m}$ và sâu $1{,}8\text{ m}$. Thể tích nước tối đa có thể chứa trong bể bơi là
\choice
{$450\text{ m}^3$}
{\True $540\text{ m}^3$}
{$600\text{ m}^3$}
{$720\text{ m}^3$}
\loigiai{
Thể tích lòng bể bơi là:
$25 \times 12 \times 1{,}8 = 540\text{ (m}^3\text{)}$.\\
\textbf{Đáp án của câu hỏi là B}
}
\end{ex}

\begin{ex} %[5D7V43]
Cửa hàng thời trang bán áo khoác với chương trình: mua chiếc thứ nhất giá $350\,000$ đồng, chiếc thứ hai được giảm giá $30\%$ so với giá chiếc thứ nhất. Bạn Minh mua 2 chiếc áo khoác đó thì phải trả tổng số tiền là
\choice
{$560\,000\text{ đồng}$}
{\True $595\,000\text{ đồng}$}
{$600\,000\text{ đồng}$}
{$630\,000\text{ đồng}$}
\loigiai{
Giá của chiếc áo khoác thứ hai là:
$350\,000 \times (100\% - 30\%) = 245\,000\text{ (đồng)}$.\\
Tổng số tiền bạn Minh phải trả là:
$350\,000 + 245\,000 = 595\,000\text{ (đồng)}$.\\
\textbf{Đáp án của câu hỏi là B}
}
\end{ex}

\begin{ex} %[5D7V43]
Một trường tiểu học tổ chức cho $420$ học sinh đi dã ngoại bằng ô tô loại $45$ chỗ ngồi. Cần thuê ít nhất bao nhiêu chiếc xe ô tô loại đó để chở hết số học sinh?
\choice
{$8\text{ xe}$}
{$9\text{ xe}$}
{\True $10\text{ xe}$}
{$11\text{ xe}$}
\loigiai{
Ta có: $420 : 45 = 9$ (dư $15$).\\
Vì còn dư $15$ học sinh nên cần thêm 1 xe nữa. Số xe ít nhất cần thuê là $9 + 1 = 10\text{ (xe)}$.\\
\textbf{Đáp án của câu hỏi là C}
}
\end{ex}

\begin{ex} %[5D7V43]
Mỗi hộp bánh quy cân nặng $0{,}45\text{ kg}$. Một thùng giấy chứa $24$ hộp bánh quy như thế. Biết vỏ thùng giấy rỗng nặng $0{,}8\text{ kg}$. Toàn bộ thùng bánh đó cân nặng số ki-lô-gam là
\choice
{$10{,}8\text{ kg}$}
{\True $11{,}6\text{ kg}$}
{$12\text{ kg}$}
{$12{,}4\text{ kg}$}
\loigiai{
Khối lượng của 24 hộp bánh quy là:
$0{,}45 \times 24 = 10{,}8\text{ (kg)}$.\\
Khối lượng của cả thùng bánh là:
$10{,}8 + 0{,}8 = 11{,}6\text{ (kg)}$.\\
\textbf{Đáp án của câu hỏi là B}
}
\end{ex}

\textbf{Phần 2. Trắc nghiệm đúng sai}
\setcounterpageref{ex}{0}

\begin{ex} %[5D7V43]
Bạn Lan mua đồ dùng học tập tại hiệu sách: $10$ quyển vở giá $9\,000$ đồng/quyển, $5$ cây bút bi giá $4\,000$ đồng/cây và $1$ bộ thước kẻ giá $15\,000$ đồng.
\choiceTF
{\True Số tiền mua vở là $90\,000$ đồng}
{\True Số tiền mua bút bi là $20\,000$ đồng}
{\True Tổng hóa đơn mua hàng của Lan là $125\,000$ đồng}
{Lan đưa cho thu ngân tờ $200\,000$ đồng thì số tiền thừa nhận lại là $85\,000$ đồng}
\loigiai{
\begin{itemchoice}
\item Đúng. $10 \times 9\,000 = 90\,000\text{ (đồng)}$.
\item Đúng. $5 \times 4\,000 = 20\,000\text{ (đồng)}$.
\item Đúng. $90\,000 + 20\,000 + 15\,000 = 125\,000\text{ (đồng)}$.
\item Sai. Số tiền thừa nhận lại là: $200\,000 - 125\,000 = 75\,000\text{ (đồng)} 
eq 85\,000\text{ đồng}$.
\end{itemchoice}
}
\end{ex}

\begin{ex} %[5D7V43]
Một căn phòng hình chữ nhật có kích thước dài $8\text{ m}$ và rộng $5\text{ m}$ được lát sàn gỗ. Tiền mua gỗ sàn là $320\,000$ đồng/$\text{m}^2$, tiền công thợ lát là $40\,000$ đồng/$\text{m}^2$.
\choiceTF
{\True Diện tích mặt sàn căn phòng là $40\text{ m}^2$}
{\True Tiền mua gỗ sàn là $12\,800\,000$ đồng}
{\True Tiền công thợ lát sàn là $1\,600\,000$ đồng}
{Tổng chi phí hoàn thiện mặt sàn căn phòng là $15\,000\,000$ đồng}
\loigiai{
\begin{itemchoice}
\item Đúng. Diện tích: $8 \times 5 = 40\text{ (m}^2\text{)}$.
\item Đúng. Tiền gỗ: $40 \times 320\,000 = 12\,800\,000\text{ (đồng)}$.
\item Đúng. Tiền công: $40 \times 40\,000 = 1\,600\,000\text{ (đồng)}$.
\item Sai. Tổng chi phí là: $12\,800\,000 + 1\,600\,000 = 14\,400\,000\text{ (đồng)} 
eq 15\,000\,000\text{ đồng}$.
\end{itemchoice}
}
\end{ex}

\textbf{Phần 3. Trắc nghiệm trả lời ngắn}
\setcounterpageref{ex}{0}

\begin{ex} %[5D7V43]
Mua $5$ quyển truyện tranh giá $24\,000$ đồng/quyển và $2$ hộp màu vẽ giá $35\,000$ đồng/hộp. Tổng số tiền phải thanh toán là bao nhiêu nghìn đồng?
\shortans{190}
\loigiai{
Tổng số tiền cần thanh toán là:
$5 \times 24\,000 + 2 \times 35\,000 = 120\,000 + 70\,000 = 190\,000\text{ (đồng)}$.\\
Đổi: $190\,000\text{ đồng} = 190\text{ nghìn đồng}$.
}
\end{ex}

\begin{ex} %[5D7V43]
Một mảnh đất hình chữ nhật có chiều dài $25\text{ m}$ và chiều rộng $18\text{ m}$. Sử dụng máy tính cầm tay tính diện tích mảnh đất theo đơn vị mét vuông.
\shortans{450}
\loigiai{
Diện tích mảnh đất là:
$25 \times 18 = 450\text{ (m}^2\text{)}$.
}
\end{ex}

\begin{ex} %[5D7V43]
Một can chứa $20\text{ l}$ dầu ăn. Mỗi lít dầu ăn cân nặng $0{,}8\text{ kg}$. Riêng vỏ can rỗng cân nặng $1\text{ kg}$. Hỏi cả can dầu đó cân nặng bao nhiêu ki-lô-gam?
\shortans{17}
\loigiai{
Khối lượng của lượng dầu ăn trong can là:
$20 \times 0{,}8 = 16\text{ (kg)}$.\\
Khối lượng cả can dầu là:
$16 + 1 = 17\text{ (kg)}$.
}
\end{ex}

\begin{ex} %[5D7V43]
Một trường tiểu học có $600$ học sinh, trong đó có $35\%$ học sinh tham gia câu lạc bộ bơi lội. Hỏi có bao nhiêu học sinh của trường tham gia câu lạc bộ bơi lội?
\shortans{210}
\loigiai{
Số học sinh tham gia câu lạc bộ bơi lội là:
$600 \times 35\% = 210\text{ (học sinh)}$.
}
\end{ex}

\begin{ex} %[5D7V43]
Mua một chiếc cặp sách giá niêm yết $250\,000$ đồng, bạn Mai được giảm giá $20\%$. Hỏi bạn Mai được giảm giá bao nhiêu nghìn đồng?
\shortans{50}
\loigiai{
Số tiền được giảm giá là:
$250\,000 \times 20\% = 50\,000\text{ (đồng)}$.\\
Đổi: $50\,000\text{ đồng} = 50\text{ nghìn đồng}$.
}
\end{ex}

\begin{ex} %[5D7V43]
Một xe ô tô chở $40$ thùng hàng, mỗi thùng hàng cân nặng $25\text{ kg}$. Tổng khối lượng số hàng mà xe ô tô đó đang chở là bao nhiêu tạ?
\shortans{10}
\loigiai{
Tổng khối lượng hàng xe chở là:
$40 \times 25 = 1\,000\text{ (kg)}$.\\
Đổi: $1\,000\text{ kg} = 10\text{ tạ}$.
}
\end{ex}
